\documentclass[12pt,a4paper]{article}

% encoding / language
\usepackage[utf8]{inputenc}
\usepackage[T1]{fontenc}
\usepackage[brazilian]{babel}

% misc
\usepackage{fmtcount}
\usepackage{amsmath,amssymb,amsthm}
\usepackage{enumitem}
\usepackage{fancyvrb}
\usepackage[dvipsnames]{xcolor}

% citations
\usepackage{csquotes}
\usepackage[style=abnt,maxbibnames=5]{biblatex}
\addbibresource{ref.bib}

% spacing
\usepackage{indentfirst}
\frenchspacing
\setlength{\parindent}{1.25cm}
\usepackage{setspace}
\onehalfspacing
\usepackage[top=3cm,bottom=2cm,left=3cm,right=2cm]{geometry}

% font
\usepackage[scaled]{helvet}
\renewcommand{\familydefault}{\sfdefault}
\usepackage{sfmath}

% hypertext
\usepackage{hyperref}
\hypersetup{
    colorlinks=true,
    urlcolor=[RGB]{0, 0, 130},
    linkcolor=.,
    citecolor=.
}
\urlstyle{same}

\newenvironment{question}[1]{%
    \pdfbookmark[3]{Questão #1}{q:#1}
    \bigskip\par\noindent
    \textbf{Questão #1} \dotfill
    \medskip\par\nopagebreak\noindent
    \ignorespaces
}{%
    \noindent\dotfill\medskip
}

\newcommand{\listtitle}[1]{%
    \medskip\par\noindent
    {\small\textbf{\underline{#1:}}}%
    \par\nopagebreak\noindent
}

\renewcommand{\theFancyVerbLine}{\ttfamily\color{gray}\arabic{FancyVerbLine}}

\title{Resumo de Complexidade de Algoritmos para Enade de Ciência da Computação}
\author{
    \small\href{https://github.com/danielsource/enade-complexidade}{Código-fonte GitHub}%
    \thanks{\url{https://github.com/danielsource/enade-complexidade}}
}
\date{\today}

\begin{document}
\maketitle
\ordinalnum{1}[f] atividade da disciplina optativa ENADE-2026.2, do Professor Ricardo Azevedo Moreira da Silva.
Este documento é destinado a estudantes de Ciência da Computação
e contém um resumo de Complexidade de Algoritmos baseado nas questões objetivas (múltipla escolha)
selecionadas das três provas de 2014, 2017 e 2021 do \citefield{enade}{title}.

\section{Contextualização}
Analisar a eficiência é o foco central da disciplina de Complexidade de Algoritmos.
Essa subárea examina o custo de tempo, memória ou outros recursos necessários para resolver problemas computacionais
(por exemplo, para decidir qual técnica utilizar dentre várias);
e a complexidade é a métrica definida de forma abstrata (em vez de usar o tempo decorrido real ou outro recurso)
para descrever a eficiência do algoritmo independentemente da sua implementação real \parencite[275--276]{sipser2012}.

Uma vez que temos um algoritmo implementado, é simples cronometrar o seu tempo de execução.
Porém, o algoritmo será mais rápido em um computador mais performático~---
que consegue executar mais instruções ou processar mais dados~---
em comparação a um computador menos performático.

Como pontuado em \textcite[27--28]{clrs2022} na análise do \textit{insertion sort},
para descobrir quanto tempo esse algoritmo de ordenação gasta, uma forma seria executá-lo em um computador.
Mas para isso, seria necessário implementar o algoritmo em uma linguagem de programação real.
O resultado desse teste demonstraria quanto tempo gastou aquela implementação, com uma certa entrada,
com algum compilador ou interpretador, usando bibliotecas específicas,
com tarefas diversas executando em plano de fundo em um sistema operacional;
esse teste é possível, porém volátil.

\section{Modelo de computação}
Dessa forma, precisamos utilizar um método que independa desses detalhes de implementação discutidos, mas nos dê informação suficiente:
analisando a ideia ou a descrição abstrata em pseudocódigo, por exemplo.
Para o propósito deste resumo, iremos nos basear numa simplificação do modelo computacional de \textcite[26]{clrs2022},
enxergando uma máquina com um processador que suporta as operações presentes no pseudocódigo e
executa cada instrução uma após a outra.

Cada execução da $k$-ésima linha de pseudocódigo tem o custo de tempo constante $c_k$.
E assim, o tempo de execução de um algoritmo será visto como a soma dos custos das linhas executadas.
Exemplificando com o pseudocódigo adaptado de \textcite[25]{clrs2022}:
\listtitle{Sum-Array(\texttt{A, n})}
\begin{Verbatim}[numbers=left, commandchars=\\\{\}]
sum = 0                 # custo = \(c\sb{1}\),    vezes = \(1\)
for i = 0 to n - 1      # custo = \(c\sb{2}\),    vezes = \(n + 1\)
    sum = sum + A[i]    # custo = \(c\sb{3}\),    vezes = \(n\)
return sum              # custo = \(c\sb{4}\),    vezes = \(1\)
\end{Verbatim}
O algoritmo \textit{Sum-Array} retorna a soma dos valores do vetor \texttt{A} de tamanho $n$.
Nesse algoritmo, temos que a linha 1 tem um custo $c_1$ e é executada $1$ vez,
já a linha 2 tem um custo $c_2$ e é executada $n + 1$ vezes, pois o intervalo do laço é fechado de $0$ até $n - 1$,
e é necessário fazer o teste da condição do laço quando $i = n$.
Iremos denotar o tempo de execução de um algoritmo com uma entrada de tamanho $n$ por $T(n)$,
obtendo isto para \textit{Sum-Array} \parencite[29]{clrs2022}:
\[
    T(n) = c_1 + c_2(n + 1) + c_3n + c_4.
\]

\section{Análise por casos}
O exemplo anterior com \textit{Sum-Array} mostra que o tempo de execução ($T(n)$) pode depender do tamanho da entrada ($n$).
Porém, existem algoritmos em que $T(n)$ também depende da própria configuração da entrada, por exemplo:
\listtitle{Linear-Search(\texttt{A, n, k})}
\begin{Verbatim}[numbers=left, commandchars=\\\{\}]
i = 0
while i < n and A[i] != k
    i = i + 1
if i == n
    return -1
return i
\end{Verbatim}
O algoritmo \textit{Linear-Search} retorna o índice do primeiro elemento do vetor \texttt{A} que seja igual a \texttt{k}, caso contrário, retorna \texttt{-1}.
Caso tenhamos que \texttt{A[0] == k}, o teste da condição do laço executará uma vez e sairá imediatamente, pois \texttt{A[i] != k} é falso.
Por outro lado, se \texttt{A} não contém \texttt{k}, o laço executará $n$ vezes e retornará \texttt{-1}.
Portanto, existe a necessidade de analisar um mesmo algoritmo em diferentes casos.

Conforme \textcite[31]{clrs2022}, geralmente nos concentramos em achar apenas o pior caso, isto é,
o maior $T(n)$ para qualquer entrada de tamanho $n$.
Ao encontrar esse pior caso, temos a garantia de que o algoritmo nunca vai demorar mais que isso.

\section{Complexidade assintótica (\texorpdfstring{$O, \Theta, \Omega$}{O, Theta, Omega})}
O tipo de análise apresentada nas seções anteriores e a presente em \textcite[seç.~2.2]{clrs2022}
ignora o custo de tempo real e, em vez disso, utiliza constantes $c_k$.
Porém, podemos ter dificuldade em analisar algoritmos cada vez maiores e com mais casos possíveis usando essa abordagem.

Como essas constantes dão mais detalhes do que o necessário, vamos eliminar todos os termos de menor grau
e todos os coeficientes da expressão $T(n)$; por exemplo, considere um algoritmo qualquer de 50 linhas de código com
\[
    T(n) = c_1 + c_2 + \dots + c_{11}8n^2 + \dots + c_{25}n^3 + \dots + c_{50}.
\]
Considere também que o termo de maior grau desse $T(n)$ é $c_{25}n^3$; logo podemos simplificar para
\[
    T(n) = n^3;
\]
sabendo que, para o nosso propósito, analisar esse novo $T(n)$ é suficiente, pois à medida que $n$ aumenta,
todos os outros termos ficam menos significativos.
Essa abstração é chamada de \emph{ordem de crescimento} (ou taxa de crescimento)
e pode ser denotada melhor usando $\Theta$ (ou \textit{Big-Theta}),
em vez de $T(n)$.
Usando o exemplo do \textit{insertion sort},
escrevemos que esse algoritmo tem o tempo de pior caso de $\Theta(n^2)$ (``theta de n ao quadrado''),
e tem o tempo de melhor caso de $\Theta(n)$.
Informalmente, entenda que $\Theta(n^2)$ significa que o custo será aproximadamente proporcional a $n^2$ quando $n$ é grande
\parencite[32--33]{clrs2022}.

Com essa métrica, é razoável, na maior parte dos casos,
considerar um algoritmo como mais eficiente em comparação a outro, se tiver uma ordem de crescimento menor.
Inevitavelmente, mesmo com grandes constantes, um algoritmo no pior caso com $\Theta(n^2)$ será mais eficiente que outro com $\Theta(n^3)$.

Assim, é possível estudar a complexidade assintótica de algoritmos.
Isto é, analisar como o custo computacional aumenta (em questão de tempo, memória, etc.) à medida que o tamanho da entrada tende para o infinito.
Enquanto a notação $\Theta$ nos informa que uma função \emph{cresce precisamente} a uma certa taxa,
a notação $O$ (ou \textit{Big-O}) nos garante que uma função \emph{não cresce mais rápido} que uma certa taxa.
O \textit{Big-O} é mais comum de ser utilizado, e também por isso, é o mais comum de ser incompreendido.
Podemos dizer que o pior caso do \textit{insertion sort} é $O(n^2)$, $O(n^3)$, ou $O(n!)$; tudo isso é verdade,
pois essa notação apenas define o limite superior.
Similarmente, a notação $\Omega$ (ou \textit{Big-Omega}) nos garante que uma função \emph{cresce pelo menos tão rápido} que uma certa taxa;
ou seja, essa notação define o limite inferior \parencite[49--51,56--57]{clrs2022}.

É importante salientar que não podemos dizer, por exemplo, que ``o algoritmo \textit{insertion sort} \emph{é} $\Theta(n^2)$'',
já que ``pior caso'' está omitido da frase; e assim, implicitamente estamos afirmando isso para todos os casos,
o que é incorreto, pois no melhor caso temos $\Theta(n)$.
Além disso, nem sempre é útil dizer que ``\textit{insertion sort} \emph{é} $O(n^2)$''; por mais que isso seja verdade,
poderíamos ter usado qualquer outro $O(n^k)$ em que $k \ge 2$ \parencite[56--57]{clrs2022}.
\newpage

\subsection{Exemplos}
Seguem agora questões resolvidas do Enade para elucidar a análise de complexidade assintótica.

\begin{question}{2014-12}
Analise o custo computacional dos algoritmos a seguir, que calculam o valor de um polinômio de grau $n$, da forma:
$P(x) = a_nx^n + a_{n-1}x^{n-1} + \dots + a_1 x + a_0$,
onde os coeficientes são números de ponto flutuante armazenados no vetor $a[0..n]$, e o valor de $n$ é maior que zero.
Todos os coeficientes podem assumir qualquer valor, exceto o coeficiente $a_n$, que é diferente de zero.

\medskip\noindent\begin{minipage}{\linewidth}
\listtitle{Algoritmo 1}
\begin{Verbatim}[numbers=left]
soma = a[0]
Repita para i = 1 até n
    Se a[i] != 0.0 então
        potência = x
        Repita para j = 2 até i
            potência = potência * x
        Fim repita
        soma = soma + a[i] * potência
    Fim se
Fim repita
Imprima(soma)
\end{Verbatim}

\listtitle{Algoritmo 2}
\begin{Verbatim}[numbers=left]
soma = a[n]
Repita para i = n-1 até 0 passo -1
    soma = soma * x + a[i]
Fim repita
Imprima(soma)
\end{Verbatim}
\medskip
\end{minipage}

\noindent
Com base nos algoritmos 1 e 2, avalie as asserções a seguir e a relação proposta entre elas.

\begin{enumerate}[label=\Roman*.]
    \item Os algoritmos possuem a mesma complexidade assintótica.
\end{enumerate}

\begin{center}
    \textbf{PORQUE}
\end{center}

\begin{enumerate}[label=\Roman*., start=2]
    \item Para o melhor caso, ambos os algoritmos possuem complexidade $O(n)$.
\end{enumerate}

\noindent
A respeito dessas asserções, assinale a opção correta.

\begin{spacing}{0.9}
\begin{small}
\begin{enumerate}[label=\Alph*)]
    \item As asserções I e II são proposições verdadeiras, e a II é uma justificativa correta da I.
    \item As asserções I e II são proposições verdadeiras, mas a II não é uma justificativa correta da I.
    \item A asserção I é uma proposição verdadeira, e a II é uma proposição falsa.
    \item A asserção I é uma proposição falsa, e a II é uma proposição verdadeira.
    \item As asserções I e II são proposições falsas.
\end{enumerate}
\end{small}
\end{spacing}
\end{question}

Antes de respondermos à questão, veremos qual é a \emph{ordem de crescimento} de cada um dos algoritmos dependendo do caso.

O Algoritmo 1 no pior caso é $\Theta(n^2)$, pois teremos uma entrada com todos os coeficientes diferentes de zero;
e dado que temos o laço externo (linhas 2--10) que executa $n$ vezes e que a condição da linha 3 é sempre verdadeira;
o laço interno (linhas 5--7) executará $i - 1$ vezes (a partir do momento em que $i >= 2$), com $i$ crescendo até $n$;
dessa forma, esse caso demonstra um laço aninhado com complexidade quadrática, ou $\Theta(n^2)$.

O Algoritmo 1 no melhor caso é $\Theta(n)$, pois teremos uma entrada com apenas um coeficiente diferente de zero ($a_n$);
e baseado no caso anterior, o laço interno (linhas 5--7) nunca será executado;
dessa forma, esse caso demonstra um laço com complexidade linear, ou $\Theta(n)$.

O Algoritmo 2 é mais eficiente e não possui múltiplos casos, já que, independentemente da entrada,
o seu laço executa $n$ vezes; portanto, é $\Theta(n)$.

Agora, podemos afirmar que a asserção I é falsa,
pois vimos que os algoritmos não possuem a mesma complexidade em todos os casos.
A asserção II é verdadeira, pois ambos os algoritmos no melhor caso possuem complexidade limitada superiormente por $n$, ou seja, são $O(n)$.
Como a asserção I é falsa, não temos a relação proposta (``asserção I porque asserção II'').
Logo, eliminando todas as opções contraditórias, temos que D é a opção correta.

\begin{question}{2014-16}
Uma pilha é uma estrutura de dados que armazena uma coleção de itens de dados relacionados e que garante o seguinte funcionamento:
o último elemento a ser inserido é o primeiro a ser removido.
É comum na literatura utilizar os nomes \textit{push} e \textit{pop} para as operações de inserção e remoção de um elemento em uma pilha, respectivamente.
O seguinte trecho de código em linguagem C define uma estrutura de dados pilha utilizando um vetor de inteiros,
bem como algumas funções para sua manipulação.

\medskip\noindent\begin{minipage}{\linewidth}
\listtitle{Pilha (corrigido)}
\begin{Verbatim}[numbers=left]
#include <stdlib.h>
#include <stdio.h>
typedef struct {
    int elementos[100];
    int topo;
} pilha;
pilha *cria_pilha() {
    pilha *p = malloc(sizeof(pilha));
    p->topo = -1;
    return p;
}
void push(pilha *p, int elemento) {
    if (p->topo >= 99)
        return;
    p->elementos[++p->topo] = elemento;
}
int pop(pilha *p) {
    int a = p->elementos[p->topo];
    p->topo--;
    return a;
}
\end{Verbatim}

O programa a seguir utiliza uma pilha.

\begin{Verbatim}[numbers=left]
int main() {
    pilha *p = cria_pilha();
    push(p, 2);
    push(p, 3);
    push(p, 4);
    pop(p);
    push(p, 2);
    int a = pop(p) + pop(p);
    push(p, a);
    a += pop(p);
    printf("%d\n", a);
    return 0;
}
\end{Verbatim}
\medskip
\end{minipage}

\noindent
A esse respeito, avalie as afirmações a seguir.

\begin{enumerate}[label=\Roman*.]
    \item A complexidade computacional de ambas as funções \texttt{push} e \texttt{pop} é $O(1)$.
    \item O valor exibido pelo programa seria o mesmo caso a instrução \texttt{a += pop(p);} fosse trocada por \texttt{a += a;}
    \item Em relação ao vazamento de memória (\textit{memory leak}), é opcional chamar a função \texttt{free(p)},
    pois o vetor usado pela pilha é alocado estaticamente.
\end{enumerate}

\noindent
É correto o que se afirma em

\begin{spacing}{0.9}
\begin{small}
\begin{enumerate}[label=\Alph*)]
    \item I, apenas.
    \item III, apenas.
    \item I e II, apenas.
    \item II e III, apenas.
    \item I, II e III.
\end{enumerate}
\end{small}
\end{spacing}
\end{question}

A afirmação I é verdadeira. As funções \texttt{push} e \texttt{pop} executam um número constante de instruções,
consistindo apenas de avaliações condicionais, atribuições e acesso indexado.
Como não possuem nenhum laço de repetição, a complexidade computacional no pior caso de ambas é $O(1)$.

A afirmação II é verdadeira. Para comprovar, podemos rastrear os valores presentes na pilha ao longo da execução das linhas da função \texttt{main}:
\begin{small}
\begin{itemize}
    \item Nas linhas 3 a 5, ocorrem os empilhamentos e a pilha fica com: \texttt{[2, 3, 4]} (com o 4 no topo).
    \item Na linha 6, \texttt{pop(p)} remove o 4, deixando a pilha assim: \texttt{[2, 3]}.
    \item Na linha 7, \texttt{push(p, 2)} adiciona o 2 no topo: \texttt{[2, 3, 2]}.
    \item Na linha 8, ocorre \texttt{int a = pop(p) + pop(p);}. Serão desempilhados o 2 e 3. O valor atribuído será $2 + 3 = 5$. A pilha fica com: \texttt{[2]}.
    \item Na linha 9, \texttt{push(p, a)} empilha o 5, resultando em: \texttt{[2, 5]}.
    \item Na linha 10, a instrução original \texttt{a += pop(p);} remove o 5 do topo e soma ao valor atual da variável \texttt{a} (que é 5),
    resultando em \texttt{a = 10}. Se a instrução fosse substituída por \texttt{a += a;}, a variável iria realizar a operação \texttt{a = 5 + 5 = 10}.
\end{itemize}
\end{small}
Desta forma, o valor final impresso pela função seria \texttt{10} em ambos os casos.

A afirmação III é falsa. Na linha 8 da implementação da pilha (\texttt{malloc(sizeof(pilha))}),
a estrutura da pilha em si é alocada dinamicamente (em uma região conhecida também como \textit{heap}),
e o vetor \texttt{elementos} faz parte desse espaço de armazenamento alocado.
Conforme a especificação da linguagem C \parencite[seç.~7.24.3]{c2y},
o tempo de vida de uma memória alocada dinamicamente (precisamente, \textit{allocated storage duration})
se estende até que ela seja explicitamente desalocada.
Por isso, a chamada à função \texttt{free(p)} não é opcional para evitar o vazamento de memória.
Como apenas as afirmações I e II estão corretas, temos que C é a opção correta.

\section{Classes de Complexidade}\label{sec:class}
Enquanto problemas como o de caminho mínimo (Dijkstra%
\footnote{Vídeo: \textit{How Dijkstra's Algorithm Works} (Canal Spanning Tree) --- \url{https://www.youtube.com/watch?v=EFg3u_E6eHU}})
e de árvore geradora mínima (Prim/Kruskal%
\footnote{Vídeo: \textit{How Do You Calculate a Minimum Spanning Tree?} (Canal Spanning Tree) --- \url{https://www.youtube.com/watch?v=Yldkh0aOEcg}})
possuem soluções eficientes em tempo polinomial $O(n^k)$ (limite superior),
um vasto conjunto de problemas de otimização é considerado intratável \parencite[cap.~21--22,34]{clrs2022}.
Em termos fundamentais, um problema é dito intratável quando não possui nenhum algoritmo exato conhecido capaz de resolvê-lo em tempo polinomial
para o pior caso, exigindo tempo de execução não-polinomial ou exponencial (por exemplo, $\Theta(2^n)$ no pior caso).
Segundo \textcite[cap.~34]{clrs2022}, temos as seguintes classes (dentre várias outras) que demonstram a diferença entre encontrar uma solução e verificar uma solução:
\begin{itemize}
    \item Classe $P$ (\textit{polynomial time}):
    Contém os problemas de decisão que podem ser \emph{resolvidos} em tempo polinomial por um algoritmo determinístico (por exemplo, em $\Theta(n^3)$ no pior caso).
    \item Classe $NP$ (\textit{non-deterministic polynomial time}):
    Contém os problemas cujas soluções candidatas (chamadas de certificados) podem ser \emph{verificadas} em tempo polinomial.
    Intuitivamente, encontrar uma solução do zero pode exigir uma busca exaustiva, mas checar se uma resposta fornecida está correta é rápido.
    Como todo problema resolvível em tempo polinomial também é verificável em tempo polinomial, temos $P \subseteq NP$.
    A principal questão em aberto na Ciência da Computação é se $P = NP$ ou $P \neq NP$.
\end{itemize}
Para mapear os problemas de forma mais precisa, definem-se os conceitos de $NP$-completo e $NP$-difícil:
\begin{itemize}
    \item $NP$-completo:\label{it:npc}
    São os problemas mais difíceis dentro da classe $NP$.
    Um problema é $NP$-completo se ele pertence a $NP$ e qualquer outro problema de $NP$ pode ser reduzido a ele em tempo polinomial.
    Intuitivamente, eles representam uma ``ponte universal'': se um único problema $NP$-completo for resolvido em tempo polinomial
    (o que ainda não se provou possível),
    todos os problemas de $NP$ também o seriam (implicando $P = NP$).
    \item $NP$-difícil (\textit{NP-hard}):
    São problemas pelo menos tão difíceis quanto os $NP$-completos, mas que não precisam pertencer à classe $NP$.
    É a categoria em que se enquadram as versões de otimização dos problemas $NP$-completos.
\end{itemize}

Um exemplo clássico da categoria $NP$ é o Problema do Caixeiro Viajante (\textit{traveling-salesperson problem}~---~TSP).
Na sua versão de decisão (``existe uma rota de custo no máximo $k$?''), o problema é $NP$-completo (e, consequentemente, $NP$-difícil),
pois basta somar as arestas da rota fornecida para verificar o certificado em tempo polinomial.
Porém, na sua versão de otimização (``qual é a rota de menor custo absoluto?''),
o problema é apenas $NP$-difícil \parencite[seç.~34.5.4]{clrs2022}.

A relevância prática dessa teoria reside no conceito de redução em tempo polinomial:
se um problema do mundo real pode ser reduzido ao TSP, ele herda sua intratabilidade.
Ao demonstrar que um problema é $NP$-difícil, pode-se interromper a busca por uma solução exata em tempo polinomial para o pior caso
e adotar abordagens de aproximação ou heurísticas capazes de encontrar soluções satisfatórias em tempo viável \parencite[1045--1047,seç.~35.2]{clrs2022}.

\subsection{Exemplos}
Seguem agora questões resolvidas do Enade para elucidar o estudo das classes de complexidade.

\begin{question}{2014-22}
Considere o processo de fabricação de um produto siderúrgico que necessita passar por $n$ tratamentos térmicos e químicos para ficar pronto.
Cada uma das $n$ etapas de tratamento é realizada uma única vez na mesma caldeira.
Além do custo próprio de cada etapa do tratamento, existe o custo de se passar de uma etapa para outra, uma vez que, dependendo da sequência escolhida,
pode ser necessário alterar a temperatura da caldeira e limpá-la para evitar a reação entre os produtos químicos utilizados.
Assuma que o processo de fabricação inicia e termina com a caldeira limpa.
Deseja-se projetar um algoritmo para indicar a sequência de tratamentos que possibilite fabricar o produto com o menor custo total.
Nessa situação, conclui-se que

\begin{spacing}{0.9}
\begin{small}
\begin{enumerate}[label=\Alph*)]
    \item a solução do problema é obtida em tempo de ordem $O(n \log n)$, utilizando-se um algoritmo ótimo de ordenação.
    \item uma heurística para a solução do problema de coloração de grafos solucionará o problema em tempo polinomial.
    \item o problema se reduz a encontrar a árvore geradora mínima para o conjunto de etapas do processo,
    requerendo tempo de ordem polinomial para ser solucionado.
    \item a utilização do algoritmo de Dijkstra para se determinar o caminho de custo mínimo entre o estado inicial e o final soluciona o problema em tempo polinomial.
    \item qualquer algoritmo conhecido para a solução do problema descrito possui ordem de complexidade de tempo não-polinomial,
    uma vez que o problema do caixeiro viajante se reduz a ele.
\end{enumerate}
\end{small}
\end{spacing}
\end{question}

O problema exige encontrar uma sequência de $n$ etapas em que cada uma é executada exatamente uma vez,
minimizando o custo total de transição e retornando ao estado inicial. Essa descrição equivale ao Problema do Caixeiro Viajante (TSP).

A opção A está incorreta.
Algoritmos de ordenação organizam elementos por chave e não resolvem problemas de menor rota em grafos com custos nas transições \parencite[cap.~8]{clrs2022}.

A opção B está incorreta.
A coloração de grafos busca atribuir cores a vértices adjacentes e não encontra circuitos de menor custo \parencite[prob.~34-3]{clrs2022}.
Além disso, heurísticas não garantem encontrar a solução ótima.

A opção C está incorreta.
A árvore geradora mínima conecta todos os vértices com o menor custo total,
mas não impõe a restrição de formar um caminho simples que visita cada etapa exatamente uma vez \parencite[cap.~21]{clrs2022}.

A opção D está incorreta.
O algoritmo de Dijkstra calcula o caminho mínimo entre dois vértices,
sem garantir que todos os demais $n$ vértices do grafo sejam visitados \parencite[cap.~22]{clrs2022}.

A opção E está correta.
Como o problema descrito se reduz ao TSP e pertence à classe $NP$-difícil na versão de otimização,
qualquer algoritmo conhecido para resolvê-lo exige tempo de execução não-polinomial no pior caso \parencite[seç.~34.5.4]{clrs2022}.

\begin{question}{2014-28}
Um cientista afirma ter encontrado uma redução polinomial de um problema $NP$-completo para um problema pertencente à classe $P$.
Considerando que esta afirmação tem implicações importantes no que diz respeito à complexidade computacional,
avalie as seguintes asserções e a relação proposta entre elas.

\begin{enumerate}[label=\Roman*.]
    \item A descoberta do cientista implica $P = NP$.
\end{enumerate}

\begin{center}
    \textbf{PORQUE}
\end{center}

\begin{enumerate}[label=\Roman*., start=2]
    \item A descoberta do cientista implica na existência de algoritmos polinomiais para todos os problemas $NP$-completos.
\end{enumerate}

\noindent
A respeito dessas asserções, assinale a opção correta.

\begin{spacing}{0.9}
\begin{small}
\begin{enumerate}[label=\Alph*)]
    \item As asserções I e II são proposições verdadeiras, e a II é uma justificativa correta da I.
    \item As asserções I e II são proposições verdadeiras, mas a II não é uma justificativa correta da I.
    \item A asserção I é uma proposição verdadeira, e a II é uma proposição falsa.
    \item A asserção I é uma proposição falsa, e a II é uma proposição verdadeira.
    \item As asserções I e II são proposições falsas.
\end{enumerate}
\end{small}
\end{spacing}
\end{question}

Conforme apresentado na seção~\ref{sec:class} (página~\pageref{it:npc}),
se um problema $NP$-completo pode ser reduzido e resolvido em tempo polinomial ($P$),
todos os demais problemas em $NP$ também terão solução polinomial (Asserção II verdadeira),
o que implica que $P = NP$ (Asserção I verdadeira).
Logo, eliminando todas as opções contraditórias, temos que A é a opção correta.

% TODO: 2014-30

% TODO: 2017-25

% TODO: 2017-33

\newpage

\section{Algoritmos de Busca e de Ordenação}

Os algoritmos de busca são utilizados para localizar um determinado elemento em uma coleção de dados.
Dependendo da estratégia utilizada, o custo computacional pode variar significativamente em função do tamanho da entrada.

Os algoritmos de ordenação, por sua vez, têm como objetivo organizar os elementos de uma coleção de dados de acordo com um determinado critério,
como a ordem crescente ou decrescente.
Diferentes estratégias podem ser utilizadas para realizar essa tarefa,
apresentando custos computacionais distintos em função do tamanho e das características da entrada.

\subsection{Exemplos}
Seguem agora questões resolvidas do Enade que envolvem os algoritmos de busca e ordenação.

\begin{question}{2021-20}
Observe o código abaixo escrito na linguagem C.

\begin{Verbatim}[numbers=left]
#include <stdio.h>
#define TAM 10
int funcao1(int vetor[], int v){
    int i;
    for (i = 0; i < TAM; i++){
        if (vetor[i] == v)
            return i;
    }
    return -1;
}
int funcao2(int vetor[], int v, int i, int f){
    int m = (i + f) / 2;
    if (v == vetor[m])
        return m;
    if (i >= f)
        return -1;
    if (v > vetor[m])
        return funcao2(vetor, v, m+1, f);
    else
        return funcao2(vetor, v, i, m-1);
}
int main(){
    int vetor[TAM] = {1, 3, 5, 7, 9, 11, 13, 15, 17, 19};
    printf("%d - %d", funcao1(vetor, 15), funcao2(vetor, 15, 0, TAM-1));
    return 0;
}
\end{Verbatim}

A respeito das funções implementadas, avalie as afirmações a seguir.

\begin{enumerate}[label=\Roman*.]
    \item O resultado da impressão na linha 24 é: 7 - 7.
    \item A função \texttt{funcao1}, no pior caso, é uma estratégia mais rápida do que a \texttt{funcao2}.
    \item A função \texttt{funcao2} implementa uma estratégia iterativa na concepção do algoritmo.
\end{enumerate}

\noindent
É correto o que se afirma em

\begin{spacing}{0.9}
\begin{small}
\begin{enumerate}[label=\Alph*)]
    \item I, apenas.
    \item III, apenas.
    \item I e II, apenas.
    \item II e III, apenas.
    \item I, II e III.
\end{enumerate}
\end{small}
\end{spacing}
\end{question}

A afirmação I é verdadeira. A \texttt{funcao1} implementa uma busca sequencial (linear),
percorrendo o vetor \texttt{\{1, 3, 5, 7, 9, 11, 13, 15, 17, 19\}} do índice \texttt{0} até encontrar o valor \texttt{15},
que está armazenado no índice \texttt{7}; logo, \texttt{funcao1(vetor, 15)} retorna \texttt{7}.
Já a \texttt{funcao2} implementa uma busca binária recursiva, iniciada com \texttt{i = 0} e \texttt{f = TAM-1 = 9}.
Na primeira chamada, \texttt{m = (0 + 9) / 2 = 4} (divisão inteira), e como \texttt{vetor[4] = 9} é menor que \texttt{15},
a função é chamada novamente com \texttt{i = 5} e \texttt{f = 9}.
Nessa segunda chamada, \texttt{m = (5 + 9) / 2 = 7}, e como \texttt{vetor[7] = 15} é igual ao valor buscado,
a função retorna \texttt{7}. Assim, a impressão na linha 24 é de fato \texttt{7 - 7}.

A afirmação II é falsa. A busca sequencial da \texttt{funcao1} possui complexidade $O(n)$ no pior caso,
pois pode ser necessário percorrer todos os \texttt{TAM} elementos do vetor até encontrar (ou não) o valor procurado.
Já a busca binária da \texttt{funcao2}, por descartar metade do espaço de busca a cada chamada recursiva,
possui complexidade $O(\log n)$ no pior caso. Portanto, no pior caso, é a \texttt{funcao2} que constitui a estratégia mais rápida, e não a \texttt{funcao1}.
A afirmação III é falsa. A \texttt{funcao2} chama a si mesma nas linhas 18 e 20 (\texttt{return funcao2(...)}),
caracterizando uma estratégia \emph{recursiva} de busca binária, e não iterativa (que exigiria o uso de um laço de repetição, como \texttt{for} ou \texttt{while}).
Como apenas a afirmação I está correta, temos que A é a opção correta.

% TODO: 2021-32

\printbibliography[heading=bibintoc]
\end{document}
